\section{Discussion \& related work}

\subsection{What was done before us}

Original work by Lambek \cite{Lambek1980} sets the tone
of the future research, establishing correspondence
between type systems based on lambda calculus with
various kinds of categories, starting with
correspondence between simply typed lambda calculus and
cartesian closed categories.

Conal Elliott's ``Compiling to Categories''
\cite{Elliott:Compcat2017} shows how to interpret (some
subset of) Haskell in arbitrary \texttt{Category}
(again, definable in Haskell) via GHC compiler plugin.
That paper is a direct inspiration for this one. There
is an important difference, though: reinterpretation is
internal to our type systems, unlocking new patterns in
the usage of reinterpretation, including (but not
limited to) universe polymorphism, external
parametricity and more. However, Conal's approach is
more flexible since his translation to categories
requires the desired category to only implement primitives
actually used for translation of a particular term at
hand.

Conor McBride's ``Crude but effective stratification''
\cite{McBride:Crude2011, McBride:Crude2012} (systemic
treatment of this idea is given by Favonia and others
\cite{Favonia:Displacement2023}) introduces an
eye-opening idea of using \textit{displacement} as a
lightweight form of universe polymorphism: to reuse term
typed at a lower universe level, lift all levels in a
term by adding a constant shift $s$. From the point of
view advocated by this paper, you can see it as an act
of reinterpreting an expression in an internal model of
a type system where the universe at level $i$ is interpreted
as a universe at level $i+s$
(Section~\ref{twoLTT-displacement}).

In ``Parametricity and Dependent Types'' by
Jean-Philippe Bernardy, Patrik Jansson and Ross Paterson
\cite{Bernardy:Parametricity2010}, authors introduce a
direct translation procedure to yield parametricity
theorems and proofs from polymorphic definitions in a
dependently typed system. The definition of <<external
parametricity>> we follow, though, is given in
``Internal Parametricity, without an Interval'' by
Thorsten Altenkirch, Yorgo Chamoun, Ambrus Kaposi and
Michael Shulman \cite{Altenkirch:Parametricity2024}.
Once again, we can reformulate the translation they give
as reinterpretation in an internal model. Altenkirch and
others noted an issue with this approach: if we apply
the external translation to the judgment with nonempty
context, context becomes parametric, too; that's the
reason they add \texttt{rel} operation and yield type
system with \textit{computational properties of higher
observational type theory}. We sidestep the issue with
context since reinterpretation is applied only to terms
whose context is also <<reinterpretable>> and is not
affected by the translation
(Section~\ref{twoLTT-parametricity}).

\subsection{What we could've done differently}

\subsubsection{Just use metaprogramming instead}

In a recent paper ``A Certifying Proof Assistant for
Synthetic Mathematics in Lean'', Wojciech Nawrocki and
others \cite{Nawrocki:Lean2026} adapt \textit{Compiling
to categories} approach \cite{Elliott:Compcat2017} to
Lean's system with dependent types, using
Categories-with-Families to represent internal models
and employing Lean's powerful metaprogramming facilities
instead of a compiler plugin. In doing so, they are able
to provide a framework for doing synthetic mathematics
in Lean, extracting proofs done in MLTT applied to more
specific categories. Thus our results in
Section~\ref{twoLTT-synthetic} are close to the results shown in
their paper. To add to that, their work is already
implemented in a practical proof assistant. Yet our type
system involves no metaprogramming since
reinterpretation is internal to it, and it handles
conversion rule. In addition, the focus is a bit
different, where we use reinterpretation to express
other polymorphic features of a language.

\subsubsection{
    Computational modality of global elements}
\label{global-elements}

Since we introduced a family of modalities of
reinterpretable terms $\Box_\ell$ in
Section~\ref{twoLTT}, one could hope to allow computation in
boxed terms. Generalizing a bit, the following
definition will be useful further:

\begin{definition}
    A dependently-typed system with a type of functions
    $(\to)$ has a \textit{computational modality of
    global elements} $\Box$ if there is a one-hole term
    $\mathrm{nec}\,\_$ and an atomic two-hole term
    $\mathrm{K}\,\_\,\_$ such that the following rules
    are admissible:
    \begin{mathpar}
    \inferrule{\cdot\vdash T\,\mathrm{type}}
        {\Gamma\vdash\Box\,T\,\mathrm{type}}
    \and
    \inferrule{\cdot\vdash t:T}
        {\Gamma\vdash\mathrm{nec}\,t:\Box\,T}
    \and
    \inferrule{
        \Gamma\vdash f:\Box\,(S\to T)\\
        \Gamma\vdash t:\Box\,S
    }{\Gamma\vdash K\,f\,t:\Box\,T}
    \end{mathpar}
    And, in addition, the following reduction is
    admissible: $K\,(\mathrm{nec}\,\lambda x.t)\,
    (\mathrm{nec}\,s)\leadsto^*
    \mathrm{nec}(t[s/x])$.
\end{definition}

With an exception of a computational rule, these are
exactly the properties of modalities $\Box_\ell$ from
the original Section~\ref{twoLTT}. The intuition behind
the reduction is that, if we apply a ``global function''
to a ``global argument'', we then should be able to
``inline'' this application, what this reduction
precisely does.

Then, a corresponding weakening of computational
reinterpretation can be defined as follows.

\begin{definition}[Modal computational reinterpretation]
    We call a type operator $F[\_]$ a modal
    computational reinterpretation of types if it is
    defined on boxed types (if
    $\Gamma\vdash\Box\,T\,\mathrm{Type}$, then
    $\Gamma\vdash F[T]\,\mathrm{Type}$), respects
    computation on types (if $T\to T'$, then
    $F[T]\to^* F[T']$) and computes as a primitive
    recursive function on normal forms of types.
    Analogously, a term with a hole
    $\llbracket\_\rrbracket$ is a modal computational
    reinterpretation of terms if it respects computation
    on terms, computes as a primitive recursive
    function on normal forms of terms and, if
    $\Gamma\vdash t:\Box\,T$, then
    $\Gamma\vdash\llbracket t\rrbracket:F[T]$.
\end{definition}

Maybe this definition of reinterpretation is ``nice''?
Regrettably, there's a no-go theorem which is similar to
na\"{\i}ve version \ref{no-go} which holds in this case:

\begin{theorem}[No-go theorem of modal reinterpretation]
    In a confluent type system with modal computational
    reinterpretation, the following should hold for
    arbitrary well-typed $s$ and $t$:
    \[
        \mathtt{app}\llbracket\mathrm{nec}\,\lambda x.t
        \rrbracket\llbracket\mathrm{nec}\,s\rrbracket
        \equiv\llbracket\mathrm{nec}(t[s/x])
        \rrbracket
    \]
    where $\equiv$ is definitional equality and
    $\mathtt{app}$ is how $\llbracket\_\rrbracket$
    acts on $K$.
\end{theorem}

\begin{proof}
    By direct example. We are given well-formed types
    $\cdot\vdash S,T\,\mathrm{type}$ and well-typed
    terms $\cdot\vdash s:S$ and $\cdot,x:S\vdash t:T$.
    Consider the following term:
    \[
        \Gamma,x:\Box(S\to T),y:\Box\,S\vdash
            \llbracket K\,x\,y\rrbracket:F[T]
    \]
    Since $K$ is atomic, the term under interpretation is in
    normal form; hence the whole term computes to
    $\mathtt{app}\llbracket x\rrbracket \llbracket y
    \rrbracket$. Substituting $\mathrm{nec}\,
    \lambda x.t$ for $x$ and $\mathrm{nec}\,s$ for $y$,
    we get $\mathtt{app}\llbracket\mathrm{nec}\,
    \lambda x.t\rrbracket\llbracket\mathrm{nec}\,
    s\rrbracket$.

    On the other hand, if we perform substitution first,
    we would get $\Gamma\vdash\llbracket K\,
    (\mathrm{nec}\,\lambda x.t)(\mathrm{nec}\,s)
    \rrbracket:F[T]$. Since modal reinterpretation
    respects reduction, this reduces to
    $\llbracket\mathrm{nec}(t[s/x])\rrbracket$.

    In a confluent type system with functions,
    multi-step reduction commutes with substitution;
    therefore, two terms are definitionally equivalent.
\end{proof}

This, of course, doesn't prohibit us from using such a
modality: as is with na\"{\i}ve interpretation, this
approach works if we manage to constrain models to those
where equalities hold definitionally.

Another benefit of modal computational reinterpretation
is the following property:

\begin{theorem}[Collapse of computational modalities]
    The following implication about $\twoLTT$ is true:
    \[\dfrac{
        \exists\kappa:
        \Gamma.\,{}_\ell\text{\faLock}\vdash_\kappa
        t:T
    }{\forall j\geqslant\mathrm{level}(T):
        \Gamma.\,{}_\ell\text{\faLock}\vdash_j
    \mathrm{nf}(t):\mathrm{nf}(T)}\]
\end{theorem}

We do not give a direct proof here, but the general idea
is that any level higher than the level of $T$ can only
appear in some redex of $t$ or $T$. Then, if modalities
were computational, it would be enough to use a single
modality $\Box$ instead of using a family of modalities
$\Box_\ell$: level of type under box would tell how many
universe models we need to provide.

But, again, it all comes down to defining the type of
\textit{definitional} internal models.

\subsubsection{Interpreting interpretation}

In Section~\ref{twoLTT}, we introduced our
``reinterpretable terms'' modality in a very
constraining way specifically so that we wouldn't have
to provide interpretation of this modality inside
internal models. Once again, the problem is with terms
that try to interpret variables from CwF context;
the simplest example is $\llbracket \lambda x.\llbracket x
\rrbracket_1\rrbracket_2$. In order to be able to
interpret such a term, we would have to define \textit{
what reinterpretation operator means inside an internal
model}. This quickly gets out of hand: rather than just
turning the type of models into a (recursive? inductive?)
type, the type of models would become such a type that a
type of one of its fields depends on the value of the
whole model. Actually, this happened before:
Kolomatskaia et al. \cite{Kolomatskaia_Shulman_2025}
have already defined a type like this inside their
\textit{displayed type theory}: type SST of
semi-simplicial types is defined coinductively in terms
of values of itself\footnote{More generally, Shulman
dubs such types \textit{modal codatatypes}
\cite{shulman2026}.}. While defining type of models in
a similar vein could be an interesting development of the type
system presented in this paper, we deem it out of scope
for now.
